\chapter{Background}\label{chap:background}
Introduce the key theoretical concepts relevant to your work.

\section{Tables}
The table below demonstrates the desired English caption: \autoref{tab:example}.
\begin{table}[htbp]
  \centering
  \caption{Example comparison of two approaches}
  \label{tab:example}
  \begin{tabular}{lcc}
    \toprule
    Criterion & Approach A & Approach B \\
    \midrule
    Simplicity & High & Medium \\
    Flexibility & Medium & High \\
    \bottomrule
  \end{tabular}
  \fonte{}
\end{table}

\section{Figures}
A figure caption uses the English label \textbf{Figure}.
\begin{figure}[htbp]
  \centering
  \fbox{\parbox[c][2.5cm][c]{0.6\textwidth}{\centering Example illustration placeholder}}
  \caption{Example illustration with an English caption}
  \label{fig:example}
  \fonte{}
\end{figure}

\section{Abbreviations}
Refer to \gls{pdf} and \gls{utfpr}. The list of abbreviations and acronyms is generated from \texttt{PreTexto/entradas-siglas.tex}.
